\section{Discussion \& Trade-Off Analysis}
\label{sec:discussion}

\subsection{Architectural Strengths}
\label{subsec:strengths}
The experimental results validate several key structural advantages of KDR-CA-AEAD:
\begin{enumerate}
    \item \textbf{Dynamic Permutation Non-Linearity}: Traditional Cellular Automata ciphers relied on fixed, un-keyed rule topologies, leaving them susceptible to algebraic reduction. By deriving rule tables dynamically per session using HKDF-SHA256, KDR-CA-AEAD eliminates static algebraic invariants.
    \item \textbf{Inter-Byte Chaining & Dual-Rule Coupling}: Candidate A-Chain state evolution couples dual rules ($R_1, R_2$) with prime offset $\delta = 13$ and inter-byte feedback ($\text{prev\_state}$). This ensures complete bit diffusion across subsequent bytes, achieving strict avalanche ratios (50.12\%) within a single forward pass.
    \item \textbf{Robust AEAD Security}: Using Encrypt-then-MAC (EtM) with HMAC-SHA256 ensures IND-CCA2 security. Modifying any single byte of ciphertext, nonce, tag, or associated data triggers instant authentication rejection.
\end{enumerate}

\subsection{Performance Trade-Offs}
\label{subsec:tradeoffs}
While software AES implementations on modern x86_64 desktop processors leverage dedicated AES-NI hardware vector units to exceed 20 MB/s in pure software, KDR-CA-AEAD provides significant advantages for lightweight edge microcontrollers and custom FPGA/ASIC hardware:
\begin{itemize}
    \item \textbf{Zero S-Box Lookup Tables}: KDR-CA-AEAD does not require pre-computed 256-byte S-boxes, mitigating cache-timing side-channel attacks.
    \item \textbf{Parallel Bit Operations}: 1D CA state updates consist of bitwise logic shifts, AND, OR, and XOR operations, requiring minimal gate counts when mapped to hardware logic cells.
\end{itemize}

\subsection{Limitations}
\label{subsec:limitations}
\begin{enumerate}
    \item \textbf{Inter-Byte Feedback Dependency}: In forward transformation, step $i$ depends on $\text{prev\_state} = T_{i-1}$. Consequently, byte-level forward encryption must be executed sequentially across a single message buffer.
    \item \textbf{Python Overhead}: High-level Python interpreter overhead limits raw single-core software throughput compared to compiled C implementations.
\end{enumerate}
